% =====================================================================
%   附录 B：操作系统 (Operating System)
% =====================================================================

\section{操作系统}\label{app:os}

\subsection{数据集细节}

\subsubsection*{构造细节}

操作系统（OS）数据集中的每条评估样本包含以下内容：
\begin{itemize}
    \item \textbf{指令（Instruction）}：以自然语言描述的、需要 LLM 解决的
          问题。
    \item \textbf{Docker 环境}：启动所使用的 Docker 镜像
          （例如预设的 \code{local-os/default}）。
    \item \textbf{初始化脚本（Initialization Script，可选）}：在交互开始之前
          需独立执行（\code{docker exec}）的 bash 脚本（例如用户配置、文件、
          系统状态）。
    \item \textbf{启动脚本（Start Script，可选）}：在 shell 创建之后、交互
          开始之前执行的 bash 脚本。
    \item \textbf{校验流水线（Checking Pipeline）}：用于判断 LLM 答案或操作
          正确性的检查方法。
    \item \textbf{示例脚本（Example Script，可选）}：作为参考解答的 bash 脚本。
          换言之，若在交互中执行该脚本，则可得到正确结果。该脚本仅用于下文
          介绍的单元测试中。
\end{itemize}

除传统的纯问答（QA）评测外，本文在 OS 评测中设计了两类任务：
\begin{itemize}
    \item \textbf{问答（QA）}：LLM 需输出命令以解决 OS 上的特定问题
          （如统计数量、查看文件内容）。在此情形下，模型最终需提交答案。
    \item \textbf{操作（Operation）}：LLM 需输出命令以在操作系统上执行
          某些可验证的操作（例如修改文件/用户状态）。在此情形下，模型无需
          提交最终答案。
\end{itemize}

借助统一的校验流水线，两类任务能够以统一的方式进行评测。

收集与 OS 相关的具有挑战性的查询是较为困难的。实践中，本文 OS 评测指令的
约一半由人工撰写或收集，另一半则主要由 \code{gpt-4} 生成，并经过严格的
单元测试筛选（即确保能产生正确答案/状态）。

对于人工编写的指令，我们首先从 Stack Overflow\footnote{\url{https://stackoverflow.com/}}
中收集了 6\,000 个带有 \code{bash} 或 \code{shell} 标签的真实问题及其解答，
并按点赞数排序。我们邀请了 8 位主修编程的标注者从中挑选具有挑战性的问题。
对每个被选中的问题，标注者需创作一条或多条任务指令，并撰写详细的问题
描述、初始化脚本、启动脚本与校验流水线。最终，我们对每条评估样本进行交叉
验证，确保其正确性。每条问题的标注平均耗时约 2 小时。

对于生成式问题，本文设计了如下单元测试环节：
\begin{enumerate}
    \item \textbf{初始化脚本正确性}：执行初始化脚本，剔除退出码不为 0
          （即初始化错误）的样本。
    \item \textbf{示例代码正确性}：执行示例代码与校验流水线，以判断答案的
          正确性。我们剔除答案错误的样本。
\end{enumerate}

最终，本文构建了 144 条高质量且多样化的 OS 评估样本，每条样本均配有
可交互的测试环境及相应的校验流水线（即脚本）。智能体以 1-shot CoT
作为提示，以获得更规范的输出格式（参见附录~\ref{app:os}~B）。

\subsubsection*{评测设置}

对每个问题（即一条指令），其执行过程可分为 3 个部分：
\begin{enumerate}
    \item \textbf{初始化（Initialization）}：根据指令创建一个指定镜像的
          Docker 容器，并运行初始化 bash 脚本以建立指令所规定的环境。
    \item \textbf{交互（Interaction）}：在该 Docker 容器中启动一个新 shell，
          并执行指令规定的启动 bash 脚本。随后将待测试的 LLM 与指令及
          问题描述一并喂入，由其开始与 shell 交互。在每一轮中，模型有
          两种可选动作：其一是运行 bash 脚本（即可在 shell 中生成并执行
          一系列命令）；其二是提交答案（即可终止交互过程）。需要注意的是，
          若超出轮数上限（默认 8 轮），则判定模型未能解决该问题。
    \item \textbf{校验（Checking）}：每个问题均对应一条包含若干脚本
          $f_1, f_2, \dots, f_n$ 的校验流水线，其中 $f_k$ 表示流水线中
          的第 $k$ 个脚本片段。对于 $f_k$，将模型的答案 $o_0$ 与此前脚本
          $f_t$（$t<k$）的输出 $o_t$ 一并作为输入参数传入，即
          $o_k = f_k(o_0, o_1, \dots, o_{k-1})$。当且仅当所有脚本的退出码
          均为 0 时，结果视为正确。
\end{enumerate}

\subsubsection*{评测指标}

本文以成功率（Success Rate）作为 LLM 在执行过程中解决问题的度量指标。
每条问题的最终状态仅有两种：正确或错误。

\subsection{动作}

在 OS 评测中，我们设计了两大类动作：
\begin{itemize}
    \item \textbf{Bash}：启动一条 bash 命令（在 \code{content} 字段中以
          文本形式输入）。
    \item \textbf{Commit}：宣布任务目标的完成。若该任务为 QA 类问题，
          则智能体应在 \code{content} 字段中提交最终答案；否则，校验流水线
          会自动检查系统状态以判断正确性。
\end{itemize}

\subsection{提示词示例}

OS 评测的提示词由指令与交互轨迹的格式化方式组成。指令提示词的一个示例如下：
\begin{quote}\small\ttfamily
You are an assistant that will act like a person, I will play the role of
linux(ubuntu) operating system. Your goal is to implement the operations
required by me or answer to the question proposed by me. For each of your
turn, you should first think what you should do, and then take exact one of
the three actions: ``bash'', ``finish'' or ``answer''.
\end{quote}
（此处译文略，下同——见原文提示词格式。详细提示词模板请参照原论文。）

轨迹采用 CoT 风格组织，并配以 1-shot 示例，以帮助模型更好地理解动作空间，
如下所示：
\begin{quote}\small\ttfamily
<USER>:\\
\hspace*{2em}\verb|{INSTRUCTION}|\\
\hspace*{2em}Now, my problem is:\\
\hspace*{2em}\verb|Tell me how many files are in the directory "/etc"?|\\
\\
<AGENT>:\\
\hspace*{2em}Think: To count the files in /etc, I need to print all the
files in it.\\
\hspace*{2em}Act: bash\\
\hspace*{4em}\verb|```bash|\\
\hspace*{4em}\verb|ls /etc|\\
\hspace*{4em}\verb|```|\\
\\
<USER>:\\
\hspace*{2em}The output of the OS:\\
\hspace*{2em}\verb|cpi cron.hourly fuse.conf iproute2 lvm ... [truncated]|\\
\\
<AGENT>:\\
\hspace*{2em}Think: The output has been truncated because it is too long,
so I need to count files by script directly.\\
\hspace*{2em}Act: bash\\
\hspace*{4em}\verb|```bash|\\
\hspace*{4em}\verb|ls -1 /etc | wc -l|\\
\hspace*{4em}\verb|```|\\
\\
<USER>:\\
\hspace*{2em}The output of the OS: 220\\
\\
<AGENT>:\\
\hspace*{2em}Think: Now I get the answer, it is 220.\\
\hspace*{2em}Act: answer(220)\\
\\
<USER>:\\
\hspace*{2em}Now, I will start a new problem in a new OS. My problem is:\\
\hspace*{2em}\verb|{PROBLEM}|
\end{quote}

LOOP 部分按 \code{<AGENT>} 回复的 \code{Action} 进行分支：若为 \code{bash}，
则执行命令并向 \code{User} 反馈输出；若为 \code{answer} 或 \code{finish}，
则执行校验脚本后退出循环；其它情形同样退出循环。
